\chapter{2010年湖南卷（文）}

\section{单选题}

\begin{problem}
复数 \(\frac{2}{1 - \i}\) 等于（\quad）

\choices
    {\(1 + \i\)}
    {\(1 - \i\)}
    {\(-1 + \i\)}
    {\(-1 - \i\)}
\end{problem}

\begin{answer}
A
\end{answer}

\begin{solution}
将分子、分母同乘分母的共轭复数，得
\[
\frac{2}{1-\i}=\frac{2(1+\i)}{(1-\i)(1+\i)}=\frac{2(1+\i)}{2}=1+\i
\]
故选 A.
\end{solution}

\begin{problem}
下列命题中的假命题是（\quad）

\choices
    {\(\exists x \in \R\)，\(\lg x = 0\)}
    {\(\exists x \in \R\)，\(\tan x = 1\)}
    {\(\forall x\in \R\)，\(x^3 >0\)}
    {\(\forall x\in \R\)，\(2^{x} > 0\)}
\end{problem}

\begin{answer}
C
\end{answer}

\begin{solution}
命题 A 中取 \(x=1\)，有 \(\lg 1=0\)，所以 A 为真；命题 B 中取 \(x=\dfrac{\pi}{4}\)，有 \(\tan\dfrac{\pi}{4}=1\)，所以 B 为真；对命题 C，取 \(x=0\)，则 \(x^3=0\) 不大于 \(0\)，所以“对任意 \(x\in\R\)，\(x^3>0\)”为假；对任意 \(x\in\R\)，\(2^x>0\)，所以 D 为真.
故选 C.
\end{solution}

\begin{problem}
某商品销售量 \( y \) （件）与销售价格 \( x \) （元/件）负相关，则其回归方程可能是（\quad）

\choices
    {\(\hat{y} = -10x + 200\)}
    {\(\hat{y} = 10x + 200\)}
    {\(\hat{y} = -10x - 200\)}
    {\(\hat{y} = 10x - 200\)}
\end{problem}

\begin{answer}
A
\end{answer}

\begin{solution}
负相关要求回归直线的斜率为负，因此排除 B、D. 回归直线必经过样本中心点 \((\bar{x},\bar{y})\)，且价格和销量均为正数，所以 \(\bar{x}>0\)、\(\bar{y}>0\). 若斜率为 \(-10\)，截距应为
\[
\bar{y}-(-10)\bar{x}=\bar{y}+10\bar{x}>0
\]
因此截距为 \(-200\) 的 C 不可能，只有 A 可能.
故选 A.
\end{solution}

\begin{problem}
极坐标 \(\rho = \cos \theta\) 和参数方程 \(\begin{cases}
x = -1 - t\\
y = 2 + t,
\end{cases}\) （\(t\) 为参数）所表示的图形分别是（\quad）

\choices
    {直线，直线}
    {直线，圆}
    {圆，圆}
    {圆，直线}
\end{problem}

\begin{answer}
D
\end{answer}

\begin{solution}
由极坐标与直角坐标的关系 \(x=\rho\cos\theta\)、\(x^2+y^2=\rho^2\)，有
\[
\rho=\cos\theta\ \Longrightarrow\ \rho^2=\rho\cos\theta\ \Longrightarrow\ x^2+y^2=x
\]
即
\[
\left(x-\frac12\right)^2+y^2=\frac14
\]
所以它表示圆. 参数方程满足
\[
x+y=(-1-t)+(2+t)=1
\]
所以它表示直线. 故选 D.
\end{solution}

\begin{problem}
设抛物线 \( y^{2} = 8x \) 上一点 \(P\) 到 \( y \) 轴的距离是 4，则点 \(P\) 到该抛物线焦点的距离是（\quad）

\choices
    {\(4\)}
    {\(6\)}
    {\(8\)}
    {\(12\)}
\end{problem}

\begin{answer}
B
\end{answer}

\begin{solution}
抛物线的标准方程为 \(y^2=2px\)，故 \(2p=8\)，从而 \(p=4\)，焦点为 \(F(2,0)\). 因为 \(P\) 到 \(y\) 轴的距离为 \(|x|=4\)，而抛物线上 \(x=\dfrac{y^2}{8}\geqslant0\)，所以 \(x=4\). 代入抛物线得 \(y^2=32\). 因此
\[
PF=\sqrt{(4-2)^2+y^2}=\sqrt{4+32}=6
\]
故选 B.
\end{solution}

\begin{problem}
若非零向量 \(\bs{a}\)，\(\bs{b}\) 满足 \(|\bs{a}| = |\bs{b}|\)，\((2\bs{a} + \bs{b}) \cdot \bs{b} = 0\)，则 \(\bs{a}\) 与 \(\bs{b}\) 的夹角为（\quad）

\choices
    {\(30^{\circ}\)}
    {\(60^{\circ}\)}
    {\(120^{\circ}\)}
    {\(150^{\circ}\)}
\end{problem}

\begin{answer}
C
\end{answer}

\begin{solution}
设 \(\bs{a}\) 与 \(\bs{b}\) 的夹角为 \(\theta\). 由题设
\[
(2\bs{a}+\bs{b})\cdot\bs{b}=0
\ \Longrightarrow\ 2\bs{a}\cdot\bs{b}+|\bs{b}|^2=0
\]
又因为 \(|\bs{a}|=|\bs{b}|\ne0\)，所以
\[
2|\bs{a}||\bs{b}|\cos\theta+|\bs{b}|^2=0
\ \Longrightarrow\ \cos\theta=-\frac12
\]
由于 \(0\leqslant\theta\leqslant\pi\)，得 \(\theta=120^\circ\). 故选 C.
\end{solution}

\begin{problem}
在 \(\triangle ABC\) 中，角 \(A\)，\(B\)，\(C\) 所对的边长分别为 \(a\)，\(b\)，\(c\)，若 \(\angle C = 120^{\circ}\)，\(c = \sqrt{2} a\)，则（\quad）

\choices
    {\(a > b\)}
    {\(a < b\)}
    {\(a = b\)}
    {\(a\) 与 \(b\) 的大小关系不能确定}
\end{problem}

\begin{answer}
A
\end{answer}

\begin{solution}
由余弦定理及 \(\cos120^\circ=-\dfrac12\)，得
\[
c^2=a^2+b^2-2ab\cos120^\circ=a^2+b^2+ab
\]
代入 \(c=\sqrt2a\)，得
\[
2a^2=a^2+b^2+ab
\ \Longrightarrow\ b^2+ab-a^2=0
\]
因 \(a>0\)、\(b>0\)，令 \(t=\dfrac ba>0\)，则
\[
t^2+t-1=0
\]
所以 \(t=\dfrac{-1+\sqrt5}{2}\)，且 \(0<t<1\). 因而 \(b<a\)，即 \(a>b\). 故选 A.
\end{solution}

\begin{problem}
函数 \( y = ax^2 + bx \) 与 \( y = \log_{\frac{|b|}{|a|}} x \)（\(ab \neq 0\)，\(|a| \neq |b|\)）在同一直角坐标系中的图象可能是（\quad）

\choices
    {\choicebitmap[width=0.15\paperwidth]{img/2010/hunan_liberal/q8_optA_clean.png}}
    {\choicebitmap[width=0.15\paperwidth]{img/2010/hunan_liberal/q8_optB_clean.png}}
    {\choicebitmap[width=0.15\paperwidth]{img/2010/hunan_liberal/q8_optC_clean.png}}
    {\choicebitmap[width=0.15\paperwidth]{img/2010/hunan_liberal/q8_optD_clean.png}}
\end{problem}

\begin{answer}
D
\end{answer}

\begin{solution}
对数函数的图象定义域为 \(x>0\)，且恒过点 \((1,0)\). 取 \(a>0\)、\(b>0\) 且 \(b<a\)，例如 \(a=2\)、\(b=1\)，则对数函数的底数
\[
\frac{|b|}{|a|}=\frac12<1
\]
所以其图象在 \(x>0\) 上单调递减并经过 \((1,0)\). 同时
\[
y=ax^2+bx=x(ax+b)
\]
是开口向上的抛物线，两个零点为 \(0\) 和 \(-\dfrac ba<0\). 这与选项 D 的图象特征相符，且满足 \(ab\ne0\)、\(|a|\ne|b|\). 故选 D.
\end{solution}

\section{填空题}

\begin{problem}
已知集合 \(A = \{1, 2, 3\}\)，\(B = \{2, m, 4\}\)，\(A \cap B = \{2, 3\}\)，则 \(m =\)\fillinblank{}.
\end{problem}

\begin{answer}
\(3\)
\end{answer}

\begin{solution}
由 \(A\cap B=\{2,3\}\)，可知 \(3\in B\). 集合 \(B=\{2,m,4\}\) 中除 \(m\) 外没有元素 \(3\)，故 \(m=3\).
\end{solution}

\begin{problem}
已知一种材料的最佳加入量在 \(100 \mathrm{~g}\) 到 \(200 \mathrm{~g}\) 之间. 若用 \(0.618\) 法安排实验，则第一次试点的加入量可以是\fillinblank{} \(\mathrm{g}\).
\end{problem}

\begin{answer}
\(161.8\)
\end{answer}

\begin{solution}
在区间 \([100,200]\) 内，按 \(0.618\) 法取靠近上端点的第一个试点，则该试点把区间按比例 \(0.618\colon 0.382\) 分割，因此加入量为 \(100+0.618(200-100)=100+61.8=161.8\). 所以填 \(161.8\).
\end{solution}

\begin{problem}
在区间 \([-1, 2]\) 上随机取一个数 \(x\)，则 \(x \in [0, 1]\) 的概率为\fillinblank{}.
\end{problem}

\begin{answer}
\(\frac{1}{3}\)
\end{answer}

\begin{solution}
在区间 \([-1,2]\) 上随机取数时，取到各长度相同子区间的可能性相同. 总区间长度为 \(2-(-1)=3\)，满足 \(x\in[0,1]\) 的区间长度为 \(1-0=1\)，所以
\[
P(x\in[0,1])=\frac{1}{3}
\]
\end{solution}

\begin{problem}
下图是求实数 \(x\) 的绝对值的算法程序框图，则判断框\(\circled{1}\)中可填\fillinblank{}.

\bitmapfigure[max width=0.40\linewidth]{img/2010/hunan_liberal/q12_fig1.png}
\end{problem}

\begin{answer}
\(x\geqslant 0\)
\end{answer}

\begin{solution}
程序在判断框中取“是”时输出 \(x\)，取“否”时输出 \(-x\). 要使输出等于 \(|x|\)，当 \(x\geqslant0\) 时应输出 \(x\)，当 \(x<0\) 时应输出 \(-x\). 因而判断框中应填 \(x\geqslant0\).
\end{solution}

\begin{problem}
图中的三个直角三角形是一个体积为 \(20 \mathrm{~cm}^{3}\) 的几何体的三视图，则 \(h=\)\fillinblank{}\(\mathrm{cm}\).

\bitmapfigure[max width=0.42\linewidth]{img/2010/hunan_liberal/q13_fig1.png}
\end{problem}

\begin{answer}
\(4\)
\end{answer}

\begin{solution}
由三个三视图可知，该几何体为三条从同一顶点出发且两两垂直的棱长分别为 \(h\)、\(5\)、\(6\) 的三棱锥. 以棱长为 \(5\)、\(6\) 的直角三角形为底，底面积为
\[
S=\frac12\times5\times6=15
\]
其高为 \(h\)，所以体积
\[
V=\frac13Sh=\frac13\times15h=5h
\]
由 \(V=20\)，得 \(5h=20\)，从而 \(h=4\).
\end{solution}

\begin{problem}
若不同两点 \(P\)，\(Q\) 的坐标分别为 \((a, b)\)，\((3 - b, 3 - a)\)，则线段 \(PQ\) 的垂直平分线 \(l\) 的斜率为\fillinblank{}，圆 \((x - 2)^2 + (y - 3)^2 = 1\) 关于直线 \(l\) 对称的圆的方程为\fillinblank{}.
\end{problem}

\begin{answer}
\(-1\)，\(x^2 + (y - 1)^2 = 1\)
\end{answer}

\begin{solution}
由
\[
Q-P=(3-b-a,3-a-b)=(3-a-b,3-a-b)
\]
因 \(P\ne Q\)，有 \(3-a-b\ne0\)，所以直线 \(PQ\) 的斜率为 \(1\)，其垂直平分线 \(l\) 的斜率为 \(-1\). 线段中点为
\[
M\left(\frac{a+3-b}{2},\frac{b+3-a}{2}\right)
\]
满足 \(x_M+y_M=3\). 因而 \(l\) 的方程为 \(x+y=3\).

设原圆圆心为 \(O(2,3)\). 点 \(O\) 到直线 \(x+y-3=0\) 的垂足为 \(H(1,2)\)，故关于 \(l\) 的对称点为
\[
O'=2H-O=(0,1)
\]
反射不改变半径，所以所求圆的方程为 \(x^2+(y-1)^2=1\).
\end{solution}

\begin{problem}
若规定 \(E = \{a_{1}, a_{2}, \dots, a_{10}\}\) 的子集 \(\{a_{i_{1}}, a_{i_{2}}, \dots, a_{i_{n}}\}\) 为 \(E\) 的第 \(k\) 个子集，其中 \(k = 2^{i_{1} - 1} + 2^{i_{2} - 1} + 2^{i_{3} - 1} + \dots + 2^{i_{n} - 1}\)，则
\begin{enumerate}
\item \(\{a_1, a_3\}\) 是 \(E\) 的第\fillinblank{} 个子集；
\item \(E\) 的第211个子集是\fillinblank{}.
\end{enumerate}
\end{problem}

\begin{answer}
\(5\)；\(\{a_{1}, a_{2}, a_{5}, a_{7}, a_{8}\}\)
\end{answer}

\begin{solution}
由题意，子集中的 \(a_i\) 对应二进制权值 \(2^{i-1}\). 因此
\[
\{a_1,a_3\}\ \text{对应的编号}=2^{1-1}+2^{3-1}=1+4=5
\]
又
\[
211=128+64+16+2+1=2^7+2^6+2^4+2^1+2^0
\]
所以第 \(211\) 个子集为
\[
\{a_8,a_7,a_5,a_2,a_1\}=\{a_1,a_2,a_5,a_7,a_8\}
\]
\end{solution}

\section{解答题}

\begin{problem}
已知函数 \( f(x) = \sin 2x - 2\sin^2 x \).
\begin{enumerate}
\item 求函数 \( f(x) \) 的最小正周期；
\item 求函数 \( f(x) \) 的最大值及 \( f(x) \) 取最大值时 \(x\) 的集合.
\end{enumerate}
\end{problem}

\begin{solution}
\begin{enumerate}
\item 由 \(2\sin ^2x=1-\cos 2x\)，得 \(f(x)=\sin 2x+\cos 2x-1=\sqrt{2}\sin\left(2x+\frac{\pi}{4}\right)-1\). 因为正弦函数的角频率为 \(2\)，所以 \(f(x)\) 的最小正周期为 \(\dfrac{2\pi}{2}=\pi\).
\item 由 \(-1\leqslant\sin\left(2x+\frac{\pi}{4}\right)\leqslant1\)，得 \(f(x)\leqslant\sqrt{2}-1\)，所以最大值为 \(\sqrt{2}-1\). 当且仅当 \(\sin\left(2x+\frac{\pi}{4}\right)=1\) 时取到最大值，即 \(2x+\frac{\pi}{4}=\frac{\pi}{2}+2k\pi\)，\(k\in\Z\)，解得 \(x=k\pi+\frac{\pi}{8}\). 因此取最大值时 \(x\) 的集合为 \(\left\{k\pi+\frac{\pi}{8}\mid k\in\Z\right\}\).
\end{enumerate}
\end{solution}

\begin{problem}
为了对某课题进行研究，用分层抽样方法从三所高校 \(A\)，\(B\)，\(C\) 的相关人员中，抽取若干人组成研究小组，有关数据见下表（单位：人）：

\begin{center}
\begin{tabular}{c|c|c}
高校 & 相关人数 & 抽取人数 \\
\hline
\(A\) & \(18\) & \(x\) \\
\(B\) & \(36\) & \(2\) \\
\(C\) & \(54\) & \(y\)
\end{tabular}
\end{center}

\begin{enumerate}
\item 求 \(x\)，\(y\).
\item 若从高校 B，C 抽取的人中选 2 人作专题发言，求这二人都来自高校 C 的概率.
\end{enumerate}
\end{problem}

\begin{solution}
\begin{enumerate}
\item 由分层抽样的等比例原则，得 \(\dfrac{x}{18}=\dfrac{2}{36}=\dfrac{y}{54}\)，所以 \(x=1\)，\(y=3\).
\item 从高校 \(B\)，\(C\) 抽取的人员中共有 \(2+3=5\) 人，任取 \(2\) 人的基本事件数为 \(\mathrm C_5^2=10\). 二人都来自高校 \(C\) 的事件数为 \(\mathrm C_3^2=3\)，故所求概率为 \(P=\dfrac{\mathrm C_3^2}{\mathrm C_5^2}=\dfrac{3}{10}\).
\end{enumerate}
\end{solution}

\begin{problem}
如图所示，在长方体 \(ABCD-A_{1}B_{1}C_{1}D_{1}\) 中，\(AB = AD = 1\)，\(AA_{1} = 2\)，\(M\) 是棱 \(CC_{1}\) 的中点.
\begin{enumerate}
\item 求异面直线 \(A_{1}M\) 和 \(C_{1}D_{1}\) 所成的角的正切值；
\item 证明：平面 \(ABM \perp\) 平面 \(A_{1}B_{1}M\).
\end{enumerate}

\bitmapfigure[max width=0.45\linewidth]{img/2010/hunan_liberal/q18_fig1.png}
\end{problem}

\begin{solution}
\begin{enumerate}
\item 建立空间直角坐标系，令 \(A(0,0,0)\)，\(B(1,0,0)\)，\(D(0,1,0)\)，\(A_1(0,0,2)\)，则 \(C(1,1,0)\)，\(C_1(1,1,2)\)，\(D_1(0,1,2)\). 因为 \(M\) 是棱 \(CC_1\) 的中点，所以 \(M(1,1,1)\). 于是 \(\overrightarrow{A_1M}=(1,1,-1)\)，\(\overrightarrow{C_1D_1}=(-1,0,0)\). 设异面直线 \(A_1M\) 和 \(C_1D_1\) 所成的锐角为 \(\theta\)，则
\[
\cos\theta=\frac{\left|\overrightarrow{A_1M}\cdot\overrightarrow{C_1D_1}\right|}{\left|\overrightarrow{A_1M}\right|\left|\overrightarrow{C_1D_1}\right|}=\frac{1}{\sqrt{3}}
\]
因此 \(\tan\theta=\sqrt{\dfrac{1}{\cos ^2\theta}-1}=\sqrt{2}\).
\item 由 \(\overrightarrow{AB}=(1,0,0)\)，\(\overrightarrow{AM}=(1,1,1)\)，平面 \(ABM\) 的一个法向量为
\[
\bs n_1=\overrightarrow{AB}\times\overrightarrow{AM}=(0,-1,1)
\]
又由 \(\overrightarrow{A_1B_1}=(1,0,0)\)，\(\overrightarrow{A_1M}=(1,1,-1)\)，平面 \(A_1B_1M\) 的一个法向量为
\[
\bs n_2=\overrightarrow{A_1B_1}\times\overrightarrow{A_1M}=(0,1,1)
\]
因为 \(\bs n_1\cdot\bs n_2=0\)，所以两平面的法向量互相垂直，从而平面 \(ABM\perp\) 平面 \(A_1B_1M\).
\end{enumerate}
\end{solution}

\begin{problem}
为了考察冰川的融化状况，一支科考队在某冰川上相距 \(8 \mathrm{~km}\) 的 \(A\)，\(B\) 两点各建一个考察基地. 视冰川面为平面形，以过 \(A\)，\(B\) 两点的直线为 \(x\) 轴，线段 \(AB\) 的垂直平分线为 \(y\) 轴建立平面直角坐标系（下图）. 考察范围到 \(A\)，\(B\) 两点的距离之和不超过 \(10 \mathrm{~km}\) 的区域.
\begin{enumerate}
\item 求考察区域边界曲线的方程；
\item 如图所示，设线段 \(P_{1}P_{2}\) 是冰川的部分边界线 （不考虑其他边界）. 当冰川融化时，边界线沿与其垂直的方向朝考察区域平行移动，第一年移动 \(0.2 \mathrm{~km}\). 以后每年移动的距离为前一年的 2 倍. 问：经过多长时间，点 \(A\) 恰好在冰川边界线上？
\end{enumerate}
\bitmapfigure[max width=0.38\linewidth]{img/2010/hunan_liberal/q19_fig1.png}
\end{problem}

\begin{solution}
\begin{enumerate}
\item 由 \(A(-4,0)\)，\(B(4,0)\)，知椭圆的两个焦点为 \(A\)，\(B\)，焦距的一半为 \(c=4\). 考察区域边界上任一点 \(P(x,y)\) 满足 \(PA+PB=10\)，所以椭圆的长半轴为 \(a=5\). 于是 \(b^2=a^2-c^2=25-16=9\)，故考察区域边界曲线的方程为 \(\dfrac{x^2}{25}+\dfrac{y^2}{9}=1\).
\item 由图中两点 \(P_1(-5,9)\)，\(P_2(-14,-3)\)，得边界线 \(P_1P_2\) 的斜率为 \(\dfrac{9-(-3)}{-5-(-14)}=\dfrac{4}{3}\)，故其方程为 \(y-9=\dfrac{4}{3}(x+5)\)，即 \(4x-3y+47=0\). 点 \(A(-4,0)\) 到该直线的距离为 \(d=\dfrac{|4\times(-4)-3\times0+47|}{\sqrt{4^2+(-3)^2}}=\dfrac{31}{5}=6.2\ \mathrm{km}\). 设经过 \(n\) 年边界线移动的总距离为 \(s_n\)，则
\[
s_n=0.2(1+2+2^2+\cdots+2^{n-1})=0.2(2^n-1)
\]
当点 \(A\) 恰好在冰川边界线上时，\(s_n=d\)，所以 \(0.2(2^n-1)=6.2\)，即 \(2^n=32\)，解得 \(n=5\). 因此经过 \(5\) 年点 \(A\) 恰好在冰川边界线上.
\end{enumerate}
\end{solution}

\begin{problem}
给出下面的数表序列：

\begin{center}
\begin{tabular}{cccc}
表 1 & 表 2 & 表 3 & \(\cdots\) \\
\(1\) & \(1\quad 3\) & \(1\)，\(3\)，\(5\) & \\
& \(4\) & \(4\quad 8\) & \\
& & \(12\) &
\end{tabular}
\end{center}

其中表 \(n\)（\(n = 1\)，\(2\)，\(3\)，\(\dots\)）有 \(n\) 行，第1行的 \(n\) 个数是 \(1\)，\(3\)，\(5\)，\(\dots\)，\(2n - 1\)，从第2行起，每行中的每个数都等于它肩上的两数之和.
\begin{enumerate}
\item 写出表 4，验证表 4 各行中数的平均数按从上到下的顺序构成等比数列，并将结论推广到表 \(n\)（\(n \geqslant 3\)） （不要求证明）；
\item 每个数列中最后一行都只有一个数，它们构成数列 1， 4， 12，…，记此数列为 \(\{b_n\}\)，求和： \(\frac{b_3}{b_1b_2} + \frac{b_4}{b_2b_3} + \cdots + \frac{b_{n+2}}{b_nb_{n+1}}\)（\(n \in \mathbf{N}^*\)）.
\end{enumerate}
\end{problem}

\begin{solution}
\begin{enumerate}
\item 表 4 为
\begin{center}
\begin{tabular}{cccc}
\(1\) & \(3\) & \(5\) & \(7\) \\
\(4\) & \(8\) & \(12\) & \\
\(12\) & \(20\) & & \\
\(32\) & & &
\end{tabular}
\end{center}
各行中数的平均数依次为 \(4\)，\(8\)，\(16\)，\(32\)，构成以 \(2\) 为公比的等比数列. 对于表 \(n\)，设第 \(r\) 行第 \(i\) 个数为 \(u_{r,i}\)，则由相邻两数相加的递推关系可得
\(u_{r,i}=2^{r-1}(2i+r-2)\)，\(1\leqslant i\leqslant n-r+1\).
所以第 \(r\) 行的平均数为
\(\frac{1}{n-r+1}\sum_{i=1}^{n-r+1}u_{r,i}=n2^{r-1}\)，\(r=1\)，\(2\)，\(\ldots\)，\(n\).
因此表 \(n\) 各行中数的平均数按从上到下的顺序构成首项为 \(n\)，公比为 \(2\) 的等比数列.
\item 最后一行只有一个数，所以由上式知 \(b_n=n2^{n-1}\). 因而
\[
\frac{b_{k+2}}{b_kb_{k+1}}=\frac{(k+2)2^{k+1}}{k2^{k-1}(k+1)2^k}=\frac{2}{k2^{k-2}}-\frac{2}{(k+1)2^{k-1}}
\]
于是所求和为
\[
\sum_{k=1}^{n}\frac{b_{k+2}}{b_kb_{k+1}}=\sum_{k=1}^{n}\left(\frac{2}{k2^{k-2}}-\frac{2}{(k+1)2^{k-1}}\right)=\frac{2}{1\cdot2^{-1}}-\frac{2}{(n+1)2^{n-1}}=4-\frac{1}{(n+1)2^{n-2}}
\]
\end{enumerate}
\end{solution}

\begin{problem}
已知函数 \( f(x) = \frac{a}{x} + x + (a - 1)\ln x + 15a \)，其中 \( a < 0 \)，且 \( a \neq -1 \).
\begin{enumerate}
\item 讨论函数 \( f(x) \) 的单调性；
\item 设函数 \( g(x) = \begin{cases}
(-2x^3 + 3ax^2 + 6ax - 4a^2 - 6a)\e^x, & x \leqslant 1,\\
\e \cdot f(x), & x > 1,
\end{cases} \) （\(\e\) 是自然数的底数）. 是否存在 \(a\)，使 \(g(x)\) 在 \([a, -a]\) 上为减函数？若存在，求 \(a\) 的取值范围；若不存在，请说明理由.
\end{enumerate}
\end{problem}

\begin{solution}
\begin{enumerate}
\item 函数 \(f(x)\) 的定义域为 \(x>0\)，且
\[
f'(x)=-\frac{a}{x^2}+1+\frac{a-1}{x}=\frac{x^2+(a-1)x-a}{x^2}=\frac{(x-1)(x+a)}{x^2}
\]
当 \(-1<a<0\) 时，\(0<-a<1\). 因为 \(x^2>0\)，所以 \(f'(x)\) 在 \((0,-a)\)、\((1,+\infty)\) 上为正，在 \((-a,1)\) 上为负. 因此 \(f(x)\) 在 \((0,-a)\)、\((1,+\infty)\) 上单调递增，在 \((-a,1)\) 上单调递减.

当 \(a<-1\) 时，\(1<-a\). 同理，\(f'(x)\) 在 \((0,1)\)、\((-a,+\infty)\) 上为正，在 \((1,-a)\) 上为负. 因此 \(f(x)\) 在 \((0,1)\)、\((-a,+\infty)\) 上单调递增，在 \((1,-a)\) 上单调递减.
\item 记
\[
P(x)=-2x^3+3ax^2+6ax-4a^2-6a
\]
当 \(x<1\) 时，\(g(x)=P(x)\e^x\)，从而
\[
g'(x)=\e^x\bigl(P(x)+P'(x)\bigr)=\e^xQ(x)
\]
其中
\[
Q(x)=-2x^3+(3a-6)x^2+12ax-4a^2
\]
又
\[
Q'(x)=-6(x+2)(x-a)
\]
若 \(a>-2\)，则 \(Q(a)=a^2(a+2)>0\)，所以 \(g'(x)\) 在 \(a\) 的右侧为正，\(g\) 不可能在 \([a,-a]\) 上为减函数. 因而必须有 \(a\leqslant-2\).

在 \(x=1\) 处，左侧函数值为 \(g(1)=\e(-4a^2+3a-2)\)，而右侧极限为 \(\lim_{x\to1+}g(x)=\e f(1)=\e(1+16a)\). 若 \(g\) 在 \([a,-a]\) 上为减函数，必须满足
\[
g(1)-\lim_{x\to1+}g(x)=-\e(4a+1)(a+3)\geqslant0
\]
结合 \(a\leqslant-2\) 时 \(4a+1<0\)，可得 \(a\geqslant-3\). 因而必要条件为 \(-3\leqslant a\leqslant-2\).

下面验证充分性. 当 \(-3\leqslant a\leqslant-2\) 时，在区间 \([a,1]\) 上有 \(x-a\geqslant0\)，且 \(Q'(x)\geqslant0\)（\(a\leqslant x\leqslant-2\)），\(Q'(x)\leqslant0\)（\(-2\leqslant x\leqslant1\)），所以 \(Q(x)\) 在 \(x=-2\) 处取得最大值. 又
\[
Q(-2)=-4(a+1)(a+2)\leqslant0
\]
故 \(Q(x)\leqslant0\)，从而 \(g'(x)\leqslant0\). 在 \(1<x\leqslant-a\) 时，\(g(x)=\e f(x)\)，且
\[
g'(x)=\e f'(x)=\e\frac{(x-1)(x+a)}{x^2}\leqslant0
\]
同时上面的函数值比较保证了 \(g\) 在 \(x=1\) 处没有向上跳跃. 因此 \(g(x)\) 在 \([a,-a]\) 上为减函数. 所以存在这样的 \(a\)，其取值范围为 \(a\in[-3,-2]\).
\end{enumerate}
\end{solution}
